\documentclass[11pt]{article}
\usepackage{ap-lit}
\title{Unit 3: Longer Fiction and Drama: Question Set}
\author{Paul Wang}
\date{November 2026}
\docpart{Unit 3 \textperiodcentered\ Question Set}

\begin{document}
\maketitle

\begin{abstract}
Two exam-style multiple-choice sets, one on a novel and one on a play, followed by a full Literary Argument essay. Passage 1 is the complete first chapter of Jane Austen's \emph{Pride and Prejudice} (1813), where an ironic narrator and a married couple's dialogue do almost all the work; Passage 2 is from Act V, Scene v of Shakespeare's \emph{Macbeth}, the ``Tomorrow, and tomorrow, and tomorrow'' speech and the scene around it. The questions concentrate on what this unit adds: irony, characterization through dialogue, and the conventions of the stage. Allow about 27 minutes for the multiple choice and 40 minutes for the essay. Paragraph numbers (novel) and line numbers (play) are printed in the margin.
\end{abstract}

\scoringline{MCQ: \blank\ / \qtotal \qquad Essay: \blank\ / 6 \qquad Target: 18+ MCQ, 4+ essay}

\section{Passage 1: Jane Austen, \emph{Pride and Prejudice}}

\begin{passage}{Questions 1--13.\quad Read the passage carefully before you choose your answers.}{The following passage is the complete first chapter of \emph{Pride and Prejudice}, a novel by Jane Austen published in 1813. Mr.\ and Mrs.\ Bennet live in the English countryside with their five unmarried daughters; Netherfield Park is a large house nearby.}
\pp It is a truth universally acknowledged, that a single man in possession of a good fortune, must be in want of a wife.
\pp However little known the feelings or views of such a man may be on his first entering a neighbourhood, this truth is so well fixed in the minds of the surrounding families, that he is considered the rightful property of some one or other of their daughters.
\pp ``My dear Mr.\ Bennet,'' said his lady to him one day, ``have you heard that Netherfield Park is let at last?''
\pp Mr.\ Bennet replied that he had not.
\pp ``But it is,'' returned she; ``for Mrs.\ Long has just been here, and she told me all about it.''
\pp Mr.\ Bennet made no answer.
\pp ``Do you not want to know who has taken it?'' cried his wife impatiently.
\pp ``\emph{You} want to tell me, and I have no objection to hearing it.''
\pp This was invitation enough.
\pp ``Why, my dear, you must know, Mrs.\ Long says that Netherfield is taken by a young man of large fortune from the north of England; that he came down on Monday in a chaise and four to see the place, and was so much delighted with it, that he agreed with Mr.\ Morris immediately; that he is to take possession before Michaelmas, and some of his servants are to be in the house by the end of next week.''
\pp ``What is his name?''
\pp ``Bingley.''
\pp ``Is he married or single?''
\pp ``Oh! Single, my dear, to be sure! A single man of large fortune; four or five thousand a year. What a fine thing for our girls!''
\pp ``How so? How can it affect them?''
\pp ``My dear Mr.\ Bennet,'' replied his wife, ``how can you be so tiresome! You must know that I am thinking of his marrying one of them.''
\pp ``Is that his design in settling here?''
\pp ``Design! Nonsense, how can you talk so! But it is very likely that he \emph{may} fall in love with one of them, and therefore you must visit him as soon as he comes.''
\pp ``I see no occasion for that. You and the girls may go, or you may send them by themselves, which perhaps will be still better, for as you are as handsome as any of them, Mr.\ Bingley may like you the best of the party.''
\pp ``My dear, you flatter me. I certainly \emph{have} had my share of beauty, but I do not pretend to be anything extraordinary now. When a woman has five grown-up daughters, she ought to give over thinking of her own beauty.''
\pp ``In such cases, a woman has not often much beauty to think of.''
\pp ``But, my dear, you must indeed go and see Mr.\ Bingley when he comes into the neighbourhood.''
\pp ``It is more than I engage for, I assure you.''
\pp ``But consider your daughters. Only think what an establishment it would be for one of them. Sir William and Lady Lucas are determined to go, merely on that account, for in general, you know, they visit no newcomers. Indeed you must go, for it will be impossible for \emph{us} to visit him if you do not.''
\pp ``You are over-scrupulous, surely. I dare say Mr.\ Bingley will be very glad to see you; and I will send a few lines by you to assure him of my hearty consent to his marrying whichever he chooses of the girls; though I must throw in a good word for my little Lizzy.''
\pp ``I desire you will do no such thing. Lizzy is not a bit better than the others; and I am sure she is not half so handsome as Jane, nor half so good-humoured as Lydia. But you are always giving \emph{her} the preference.''
\pp ``They have none of them much to recommend them,'' replied he; ``they are all silly and ignorant like other girls; but Lizzy has something more of quickness than her sisters.''
\pp ``Mr.\ Bennet, how \emph{can} you abuse your own children in such a way? You take delight in vexing me. You have no compassion for my poor nerves.''
\pp ``You mistake me, my dear. I have a high respect for your nerves. They are my old friends. I have heard you mention them with consideration these last twenty years at least.''
\pp ``Ah, you do not know what I suffer.''
\pp ``But I hope you will get over it, and live to see many young men of four thousand a year come into the neighbourhood.''
\pp ``It will be no use to us, if twenty such should come, since you will not visit them.''
\pp ``Depend upon it, my dear, that when there are twenty, I will visit them all.''
\pp Mr.\ Bennet was so odd a mixture of quick parts, sarcastic humour, reserve, and caprice, that the experience of three-and-twenty years had been insufficient to make his wife understand his character. \emph{Her} mind was less difficult to develop. She was a woman of mean understanding, little information, and uncertain temper. When she was discontented, she fancied herself nervous. The business of her life was to get her daughters married; its solace was visiting and news.
\end{passage}

\begin{q}{3}{3.4 Ironic narrator}{Paragraph 2 reveals that the ``want'' belongs to the neighbors with daughters, not the rich man.}
The first sentence of the novel (paragraph 1) is best read as
\ch{a sincere statement of a law of human nature that the novel will prove}{Mr.\ Bennet's opinion, reported by the narrator}{ironic: the ``truth'' is really the conviction of families eager to marry off their daughters}{a summary of Mr.\ Bingley's own reasons for coming to Netherfield}
\end{q}

\begin{q}{1}{3.5 Word choice}{``Rightful property'' makes the rich man a possession to be claimed, inverting the first sentence.}
In paragraph 2, the phrase ``the rightful property of some one or other of their daughters'' chiefly suggests that
\ch{the neighborhood regards a wealthy newcomer as a possession to be claimed}{Mr.\ Bingley has already chosen a wife}{the daughters own land in the neighborhood}{the families are indifferent to marriage}
\end{q}

\begin{q}{4}{3.1 Characterization}{His one-line, reported replies withhold what she wants: a teasing reticence against her eagerness.}
Paragraphs 4 and 6 (``Mr.\ Bennet replied that he had not''; ``Mr.\ Bennet made no answer'') chiefly serve to
\ch{show that Mr.\ Bennet has not heard his wife}{suggest that Mr.\ Bennet is ill}{indicate that the narrator does not know what Mr.\ Bennet said}{set Mr.\ Bennet's deliberate reticence against his wife's eagerness to talk}
\end{q}

\begin{q}{2}{3.5 Wit}{He names her real motive (she wants to tell), with dry, knowing wit.}
Mr.\ Bennet's reply in paragraph 8 (``\emph{You} want to tell me, and I have no objection to hearing it'') reveals chiefly his
\ch{eagerness to learn about the new tenant}{dry wit in exposing that the conversation serves his wife's wish, not his}{anger at Mrs.\ Long for gossiping}{confusion about who has taken Netherfield}
\end{q}

\begin{q}{2}{3.4 Narrator's comment}{``Invitation enough'': the narrator wryly notes how little encouragement she needs.}
The narrator's comment ``This was invitation enough'' (paragraph 9) is best described as
\ch{a neutral report of a formal invitation}{a wry observation that Mrs.\ Bennet needs little encouragement to talk}{evidence that Mr.\ Bennet has invited Mr.\ Bingley to dinner}{a shift into Mrs.\ Bennet's point of view}
\end{q}

\begin{q}{3}{3.1 Characterization through syntax}{A breathless chain of ``that'' clauses piles up secondhand detail: an eager gossip.}
The structure of Mrs.\ Bennet's speech in paragraph 10, a single long sentence of clauses beginning ``that he \ldots,'' characterizes her as
\ch{precise and well informed}{reluctant to share what she has heard}{an eager gossip piling up secondhand detail}{suspicious of the new tenant}
\end{q}

\begin{q}{1}{3.5 Feigned naivety}{Pretending to think the plan is Bingley's exposes that the ``design'' is hers.}
Mr.\ Bennet's question ``Is that his design in settling here?'' (paragraph 17) is best understood as
\ch{mock innocence that exposes the matchmaking ``design'' as his wife's, not Bingley's}{a sincere request for information about Bingley's plans}{evidence that Mr.\ Bennet suspects Bingley of a crime}{a sign that he has already met Bingley}
\end{q}

\begin{q}{4}{3.1 Character}{She takes his teasing as a real compliment: vanity and a literal mind.}
Mrs.\ Bennet's reply in paragraph 20 (``My dear, you flatter me'') reveals that she
\ch{sees through her husband's teasing and answers it in kind}{is offended by her husband's remark}{has decided not to visit Mr.\ Bingley}{takes her husband's teasing as a sincere compliment}
\end{q}

\begin{q}{3}{3.5 Verbal irony}{Calling her nerves his ``old friends'' mocks twenty years of complaint.}
In paragraph 29, Mr.\ Bennet's claim that his wife's nerves ``are my old friends'' is an example of
\ch{dramatic irony, since Mrs.\ Bennet knows something he does not}{a sincere expression of sympathy}{verbal irony mocking her habit of complaining about her nerves}{an allusion to a medical text}
\end{q}

\begin{q}{2}{3.2 Social setting}{``Four or five thousand a year,'' ``an establishment'': marriage as economic settlement.}
Mrs.\ Bennet's emphasis on Bingley's ``four or five thousand a year'' (paragraph 14) and on ``what an establishment it would be'' for a daughter (paragraph 24) chiefly reveals
\ch{her ignorance of money}{a social world in which marriage is above all an economic settlement}{that the Bennets are wealthier than their neighbors}{her wish that her daughters remain unmarried}
\end{q}

\begin{q}{4}{3.3 Structure}{After pages of dialogue, the narrator's summary confirms and sharpens what the talk has shown.}
The final paragraph (paragraph 34) differs from most of the chapter in that it
\ch{introduces Mr.\ Bingley for the first time}{reports another conversation with Mrs.\ Long}{reveals a plot twist that the dialogue concealed}{uses direct characterization to confirm and sharpen what the dialogue has already shown}
\end{q}

\begin{q}{1}{3.4 Narrator's tone}{``Mean understanding, little information'': a cutting, ironic judgment of Mrs.\ Bennet.}
The narrator's description of Mrs.\ Bennet in paragraph 34 (``a woman of mean understanding, little information, and uncertain temper'') is best described as
\ch{cuttingly ironic}{affectionately admiring}{deliberately neutral}{fearful}
\end{q}

\begin{q}{3}{3.6 Literary argument}{Twenty-three years without understanding: wit and urgency that never meet.}
Which claim about the chapter is best supported?
\ch{Mr.\ and Mrs.\ Bennet share the same goals but disagree only about manners.}{Mrs.\ Bennet is the more perceptive of the two.}{The chapter presents the Bennets' marriage as a mismatch in which the husband's wit and the wife's urgency never truly meet.}{Mr.\ Bennet refuses ever to visit Mr.\ Bingley.}
\end{q}

\section{Passage 2: William Shakespeare, \emph{Macbeth}}

\begin{poem}{Questions 14--25.\quad Read the passage carefully before you choose your answers.}{The following passage is from Act V, Scene v of \emph{Macbeth} (first performed around 1606; published 1623). Macbeth, who seized the Scottish crown by murdering King Duncan, waits in his castle at Dunsinane as an army led by Duncan's son Malcolm approaches. Witches earlier told him that he would never be defeated until Birnam Wood came to Dunsinane. Seyton is his officer.}
\sd{Enter, with drum and colours, Macbeth, Seyton, and Soldiers.}
\spk{MACBETH}
\pl{Hang out our banners on the outward walls;}
\pl{The cry is still, ``They come'': our castle's strength}
\pl{Will laugh a siege to scorn: here let them lie}
\pl{Till famine and the ague eat them up:}
\pl{Were they not forc'd with those that should be ours,}
\pl{We might have met them dareful, beard to beard,}
\pl{And beat them backward home.}
\sd{A cry of women within.}
\pl{\hspace*{8em}What is that noise?}
\spk{SEYTON}
\pl{It is the cry of women, my good lord.}
\sd{Exit.}
\spk{MACBETH}
\pl{I have almost forgot the taste of fears:}
\pl{The time has been, my senses would have cool'd}
\pl{To hear a night-shriek; and my fell of hair}
\pl{Would at a dismal treatise rouse and stir}
\pl{As life were in't: I have supp'd full with horrors;}
\pl{Direness, familiar to my slaught'rous thoughts,}
\pl{Cannot once start me.}
\sd{Re-enter Seyton.}
\pl{\hspace*{8em}Wherefore was that cry?}
\spk{SEYTON}
\pl{The queen, my lord, is dead.}
\spk{MACBETH}
\pl{She should have died hereafter;}
\pl{There would have been a time for such a word.}
\pl{To-morrow, and to-morrow, and to-morrow,}
\pl{Creeps in this petty pace from day to day,}
\pl{To the last syllable of recorded time;}
\pl{And all our yesterdays have lighted fools}
\pl{The way to dusty death. Out, out, brief candle!}
\pl{Life's but a walking shadow; a poor player,}
\pl{That struts and frets his hour upon the stage,}
\pl{And then is heard no more: it is a tale}
\pl{Told by an idiot, full of sound and fury,}
\pl{Signifying nothing.}
\sd{Enter a Messenger.}
\pl{Thou com'st to use thy tongue; thy story quickly.}
\spk{MESSENGER}
\pl{Gracious my lord,}
\pl{I should report that which I say I saw,}
\pl{But know not how to do it.}
\spk{MACBETH}
\pl{\hspace*{8em}Well, say, sir.}
\spk{MESSENGER}
\pl{As I did stand my watch upon the hill,}
\pl{I look'd toward Birnam, and anon, methought,}
\pl{The wood began to move.}
\spk{MACBETH}
\pl{\hspace*{8em}Liar, and slave!}
\end{poem}

\begin{q}{4}{3.1 Character}{Banners, a castle that ``laughs'' at a siege: defiant bravado, with resentment at deserters.}
In lines 1--7, Macbeth's attitude toward the approaching army is best described as
\ch{fearful and ready to surrender}{indifferent to the outcome}{grief-stricken over his losses}{defiant and contemptuous, with resentment at the soldiers who have deserted him}
\end{q}

\begin{q}{1}{3.5 Personification}{A castle that will ``laugh a siege to scorn'' voices Macbeth's own contempt.}
The personification in lines 2--3 (``our castle's strength / Will laugh a siege to scorn'') chiefly conveys
\ch{Macbeth's contemptuous confidence in his defenses}{the soldiers' fear of the siege}{the castle's age and weakness}{Macbeth's wish to make peace}
\end{q}

\begin{q}{2}{3.2 Situation}{``Forc'd with those that should be ours'': his own men have joined the enemy.}
Line 5 (``Were they not forc'd with those that should be ours'') reveals that
\ch{Malcolm's army is smaller than Macbeth's}{some of Macbeth's own soldiers have deserted to the enemy}{Macbeth plans to attack at once}{the witches have abandoned Macbeth}
\end{q}

\begin{q}{4}{3.4 Self-revelation}{He was once terrified by a night-shriek; now horror cannot ``start'' him.}
In lines 10--16, Macbeth reflects chiefly that
\ch{he regrets killing Duncan}{he is afraid of the cry he has just heard}{the women of the castle are in danger}{he has grown so hardened by horror that he can no longer feel fear}
\end{q}

\begin{q}{3}{3.5 Metaphor}{He has eaten his fill of horrors: sated, numbed, past feeling.}
The metaphor ``I have supp'd full with horrors'' (line 14) presents horror as
\ch{a disease he has caught}{a story he has told}{a meal he has eaten so fully that he can no longer taste it}{a battle he has won}
\end{q}

\begin{q}{1}{3.4 Ambiguity}{A muted response: she would have died later anyway, or there is no time now to mourn.}
Macbeth's response to the queen's death, ``She should have died hereafter; / There would have been a time for such a word'' (lines 19--20), is best read as
\ch{a numbed reply suggesting there is no fitting time for grief now}{an expression of relief that she is gone}{a vow to avenge her death on Malcolm}{an order to bury her with honor}
\end{q}

\begin{q}{2}{3.5 Repetition}{``To-morrow'' three times, then ``creeps \ldots\ petty pace'': days as a slow, meaningless march.}
The repetition of ``To-morrow'' in line 21, followed by ``Creeps in this petty pace from day to day,'' chiefly conveys
\ch{hope for a better future}{the slow, meaningless succession of days}{Macbeth's plans for the coming battle}{the passing of seasons on the heath}
\end{q}

\begin{q}{4}{3.5 Metaphor}{Each past day has been a light that only guided fools toward the grave.}
In lines 24--25, ``all our yesterdays have lighted fools / The way to dusty death'' suggests that
\ch{the past offers wisdom for the future}{death can be avoided by the wise}{Macbeth blames his wife for his crimes}{every day lived has only lit the way toward death}
\end{q}

\begin{q}{3}{3.5 Extended metaphor}{Life as a ``poor player'' who ``struts and frets'' and is then forgotten.}
Lines 26--28 compare life to
\ch{a candle that never goes out}{a long journey with a clear destination}{an actor's brief, self-important performance, forgotten afterward}{a king's proud procession}
\end{q}

\begin{q}{1}{3.4 Metatheatre}{An actor on a stage calls life a ``poor player'': the audience is made aware of the play itself.}
Because the lines about ``a poor player, / That struts and frets his hour upon the stage'' are spoken by an actor on a stage, the comparison is
\ch{self-referential, reminding the audience that they are watching just such a player}{inaccurate, since Macbeth is a king and not an actor}{a stage direction that the actor reads aloud}{an aside to Seyton that the audience is not meant to hear}
\end{q}

\begin{q}{2}{3.3 Short line}{The half-line falls off into silence, enacting the ``nothing'' it names.}
The speech ends with the short, incomplete line ``Signifying nothing'' (line 30). The brevity of the line chiefly
\ch{signals that Macbeth has been interrupted by the messenger}{lands with abrupt finality, leaving a silence that enacts the ``nothing'' it names}{shows that the speech is written in prose}{introduces a rhyme that ends the scene}
\end{q}

\begin{q}{4}{3.3 Dramatic irony}{The prophecy starts to come true, and numb despair snaps into rage.}
The messenger's report and Macbeth's response in lines 36--39 (``The wood began to move.'' ``Liar, and slave!'') chiefly serve to
\ch{provide comic relief before the battle}{show that Macbeth has accepted his fate calmly}{reveal that the messenger is lying}{jolt Macbeth from despair into rage as the prophecy he trusted begins to turn against him}
\end{q}

\section{Free Response: Literary Argument}

\begin{frq}{Question 3 \textperiodcentered\ Literary Argument \textperiodcentered\ 6 points \textperiodcentered\ 40 minutes}
Many works of literature feature a character who keeps up an appearance that conceals a different reality. The distance between the role a character plays in public and the self beneath it can reveal a great deal about that character and about the world of the work.

Either from your own reading or from the list below, choose a work of fiction in which a character maintains an appearance that hides something important about their nature or situation. Then, in a well-written essay, analyze how the character's effort to maintain that appearance contributes to an interpretation of the work as a whole. Do not merely summarize the plot.

In your response you should do the following:
\begin{itemize}[leftmargin=1.4em]
\item Respond to the prompt with a thesis that presents a defensible interpretation.
\item Select and use evidence to support your line of reasoning.
\item Explain how the evidence supports your line of reasoning.
\item Use appropriate grammar and punctuation in communicating your argument.
\end{itemize}

\medskip
\begin{multicols}{3}\raggedright\footnotesize
\emph{The Age of Innocence}\\
\emph{The Awakening}\\
\emph{A Doll's House}\\
\emph{Death of a Salesman}\\
\emph{Great Expectations}\\
\emph{The Great Gatsby}\\
\emph{Hamlet}\\
\emph{The Importance of Being Earnest}\\
\emph{Invisible Man}\\
\emph{Jane Eyre}\\
\emph{Macbeth}\\
\emph{Passing}\\
\emph{The Picture of Dorian Gray}\\
\emph{Pride and Prejudice}\\
\emph{The Remains of the Day}\\
\emph{The Scarlet Letter}\\
\emph{Their Eyes Were Watching God}\\
\emph{Twelfth Night}
\end{multicols}
\end{frq}

\makekey

\subsection*{Essay scoring guide (6 points)}
\small
\begin{compare}{Row}{Earns the point(s)}{Does not earn the point(s)}
A. Thesis (1) & Names a work and a character, identifies the appearance and what it hides, and makes a defensible claim about how maintaining it shapes the meaning of the \emph{whole} work. & ``In \emph{The Scarlet Letter}, Dimmesdale hides his sin''; a thesis about the character with no claim about the work; a statement that appearances can be deceiving. \\
B. Evidence and commentary (4) & \textbf{4}: specific, accurate references to several moments across the work, including its ending, with commentary that consistently explains how the character's effort (its causes, its costs, its collapse or success) builds the interpretation. \textbf{3}: specific evidence supporting a line of reasoning, with some of it explained. \textbf{2}: some specific evidence, but explanation that is general, uneven, or slides into summary. \textbf{1}: mostly plot summary or vague references. & Retelling the plot in order; references that are inaccurate; evidence from only the opening chapters. \\
C. Sophistication (1) & A complex argument across the essay, e.g., that the appearance is both a protection and a prison; that the work judges the society that requires the mask as much as the character who wears it; or that the ending's revelation is itself ambiguous. & A final sentence about how ``we all wear masks''; an unconnected reference to the author's life. \\
\end{compare}

\begin{modelessay}[A 6-point response on \emph{The Scarlet Letter}, in outline]
\textbf{Thesis.} In Nathaniel Hawthorne's \emph{The Scarlet Letter} (1850), the Reverend Arthur Dimmesdale's effort to preserve his reputation for holiness while concealing that he is the father of Hester Prynne's child becomes the novel's central argument against Puritan hypocrisy: the public sin Hester is forced to wear slowly becomes a source of strength, while the private sin Dimmesdale hides consumes him, so that the novel finds more cruelty in concealment than in shame.

\textbf{Body 1: the mask and its cost.} On the scaffold in the opening chapters, Dimmesdale urges Hester to name her lover, a speech whose eloquence is itself his disguise. His congregation's reverence grows as his guilt deepens; the more he confesses in vague terms that he is ``vile,'' the holier they think him. The appearance is not only maintained but fed by the very guilt it hides.

\textbf{Body 2: the body as evidence.} Dimmesdale's hand pressed over his heart, his wasting health, and his secret penance show the concealment written on his body, and Chillingworth, the physician who becomes his tormentor, reads it there. The midnight vigil on the scaffold, where he stands with Hester and Pearl only in darkness, shows his wish to confess and his refusal to do it where it would count.

\textbf{Body 3: the ending.} After the Election Sermon, his greatest public triumph, Dimmesdale finally mounts the scaffold in daylight and confesses, then dies. Hawthorne's narrator refuses to settle what the crowd saw on his breast, but draws the moral directly: ``Be true! Be true! Be true!'' The ending completes the argument, and its ambiguity shows how a community built on appearances cannot agree even on what it has seen.

\textbf{Why it earns Row C.} The essay argues that the novel indicts the society that rewards the mask as much as the man who wears it, and it treats the ending's ambiguity as part of the meaning rather than a problem.
\end{modelessay}

\end{document}
